\section{Results}
\label{sec:results}

% NOTE: This section is a structural skeleton.  Every \todopilot{...}
% marker is a place where pilot data needs to be inserted once the
% nine training runs finish.  See docs/research_playbook.md for the
% commands that produce these data; figures live in paper/figures/.

\subsection{Setup recap}

Three observation treatments---\emph{pixel}, \emph{symbolic},
\emph{hybrid}---trained for $10^7$ environment steps each, three seeds
per treatment, identical hyperparameters
(\Cref{sec:hyperparameters}).  All other components of the training
pipeline are matched (\Cref{sec:methods}).

\subsection{Learning Curves}
\label{sec:results-curves}

\begin{figure}[t]
\centering
% Produced by scripts/analyze.py and scripts/compare.py
% \includegraphics[width=0.85\linewidth]{figures/learning_curves.pdf}
\caption{%
  \todopilot{Learning curves: per-episode total reward smoothed with a
    20-episode rolling mean, with mean ($\pm$ std) bands across the
    three seeds for each treatment.  Replace this caption with one or
    two sentences describing the qualitative shape: which treatment
    rises fastest, which plateaus, whether bands overlap.}
}
\label{fig:learning-curves}
\end{figure}

\todopilot{Insert one paragraph describing \Cref{fig:learning-curves}
qualitatively.  Mention: the treatment with the steepest rise, the
treatment with the highest asymptote within the 10M-step budget,
whether the curves cross, and whether the across-seed std bands
overlap meaningfully.}

\subsection{Aggregate Metrics}
\label{sec:results-aggregate}

\begin{table}[t]
\centering
\small
\caption{%
  Aggregate metrics per treatment.  IQM is the interquartile mean of
  per-seed best-episode reward; CIs are 95\% percentile bootstrap
  intervals with 2{,}000 resamples.  $n$ = 3 seeds per treatment.
}
\label{tab:aggregate}
\begin{tabular}{lcccc}
\toprule
Treatment & Mean & Median & IQM & 95\% CI \\
\midrule
Pixel    & \todopilot{x.x} & \todopilot{x.x} & \todopilot{x.x} & \todopilot{[lo, hi]} \\
Symbolic & \todopilot{x.x} & \todopilot{x.x} & \todopilot{x.x} & \todopilot{[lo, hi]} \\
Hybrid   & \todopilot{x.x} & \todopilot{x.x} & \todopilot{x.x} & \todopilot{[lo, hi]} \\
\bottomrule
\end{tabular}
\end{table}

\todopilot{Insert one paragraph interpreting \Cref{tab:aggregate}.
Highlight: the highest-IQM treatment, whether its CI overlaps the
others, and the relative spread of the bands across seeds (e.g., "the
hybrid treatment had a tighter CI, suggesting more consistent
training across seeds").}

\paragraph{Probability of improvement.}
\todopilot{Insert pairwise PoI numbers from
\texttt{scripts/analyze.py}'s probability-of-improvement matrix.  For
each ordered pair (A, B), report $\Pr[\text{score}_A >
\text{score}_B]$ with its 95\% CI.  Mention which pairs cross or fail
to cross 0.5.}

\subsection{Milestone Progression}
\label{sec:results-milestones}

\begin{figure}[t]
\centering
% \includegraphics[width=0.95\linewidth]{figures/milestone_race.pdf}
\caption{%
  \todopilot{First-episode race chart for the 18 pre-registered event
    flags.  For each treatment, the median first-triggered episode
    across seeds is plotted; flags never triggered by any seed are
    shown as a hatched bar at the right edge.  Lower is faster.}
}
\label{fig:milestones}
\end{figure}

\todopilot{Insert one paragraph reporting which milestones each
treatment reached.  Pay particular attention to: (i) the
\texttt{EVENT\_GOT\_STARTER} flag, which any minimally-functional
agent should reach quickly; (ii) the \texttt{EVENT\_BEAT\_BROCK} flag,
which is the primary metric and is unlikely to be reached at 10M
steps; (iii) the canary flag
\texttt{EVENT\_GOT\_POKEMON\_FROM\_FAN\_CLUB\_CHAIRMAN}, which should
\emph{not} be reached---confirming that the agent is not exploiting an
unintended path.}

\subsection{Final-Window Performance}
\label{sec:results-final}

\begin{figure}[t]
\centering
% \includegraphics[width=0.6\linewidth]{figures/final_performance.pdf}
\caption{%
  Final-window performance: mean reward over the last 50 episodes per
  seed, aggregated as mean $\pm$ std across the three seeds per
  treatment.
}
\label{fig:final-perf}
\end{figure}

\todopilot{One short paragraph: agreement or disagreement between the
final-window bar chart and the IQM table.  Disagreement would be
informative---e.g., a treatment with high IQM but low final-window
reward suggests instability late in training.}

\subsection{Compute}
\label{sec:results-compute}

\todopilot{Add a small table with wall-clock per run (mean across
seeds), GPU model, peak GPU utilisation, and total GPU-hours for the
nine pilot runs.  Pull from \texttt{paper/compute\_ledger.md}.}
